En el ámbito de la inteligencia artificial la separación de fuentes musicales es un campo que ha experimentado gran evolución desde sus primeros acercamientos en 1973, donde se utilizaba un \textit{vocoder homomórfico}\cite{rafii18} para separar la función de excitación de la respuesta al impulso del tracto vocal en voces. A lo largo de su evolución, los métodos de separación de fuentes musicales se han desarrollado a partir de distintos paradigmas metodológicos, entre los que destacan el procesamiento de señales, el modelado o la modulación del audio y del discurso, y los enfoques probabilísticos o guiados por datos. Esta clasificación general servirá como marco para estructurar el presente capítulo.

Primeramente en el grupo de procesado de señales se persigue obtener un filtrado de la mezcla. Esto se logra comúnmente aplicando un filtro de pasa baja, pasa alta o de banda a la mezcla completa. También engloba esta rama la generación de máscaras de ganancia, que son máscaras de valores en rango [0-1] a aplicar sobre una mezcla y que devuelven una cadena de valores correspondientes a la representación tiempo-frecuencia de la fuente buscada.

En segundo lugar, en la modelación de audio o modulación de discurso, se estudia la predictibilidad del comportamiento de las señales haciendo diferenciación entre \textit{señales sinusoidales} o tonos puros y \textit{ruido}, siendo ambos los extremos en una línea de predictibilidad. Por un lado el comportamiento de las señales sinusoidales puede ser predicho en un futuro a partir de muestreos previos, mientras que por otro, el \textit{ruido blanco}, que se define como una señal impredecible a futuro y su espectrograma presenta energía constante en todo el rango frecuencial. Dentro de esta categoría se presentan varias herramientas como por ejemplo \textit{cepstrum} \cite{cepstrum}, con aplicaciones de estimación de frecuencias fundamentales, derivación de frecuencias cepstrales, y filtrado de señal.

En los acercamientos probabilísticos, se diferencia el \textit{Modelo oculto de Markov}, modelo específico de una red bayesiana que precisamente tuvo su primer uso en reconocimiento de habla \cite{hiddenmarkov} y el \textit{modelo de mezcla Gaussiano} que da una forma de aproximar una distribución como una suma ponderada de \textit{gaussianos} \cite{gaussianMixModels}, útil para crear máscaras.